\subsection*{控制流：发散、并发与协同}

每一种面向并行机的语言都必须回答同一个问题：\textbf{当不同线程需要做不同的事情时，会发生什么？}

CPU 的各个核心具有独立的程序计数器；共享指令流的线程遇到发散路径时，则需要依次执行各条路径，并屏蔽当前不执行的线程。

所有线程都求值该条件，条件为假的线程被屏蔽（停用）；这就是\textbf{谓词执行}——发散路径可能需要额外执行时间。

分块矩阵乘法还包含职责不同的工作：一个组装入下一批操作数，另一个组使用当前操作数计算。两组需要在明确的交接点之间独立推进。运行时条件描述实例选择哪个分支，而流水线还需要表达这些独立职责及其依赖关系。

Croqtile 引入两个额外的控制流原语来填补这个空白：

\begin{itemize}[leftmargin=1.6em]
\item \textbf{\texttt{inthreads.async}} —— \textbf{结构化并发区域}：在编译时把线程分成若干组，各组\textbf{同时运行不同的程序}，编译器为它们生成彼此独立的指令流；
\item \textbf{\texttt{shared event} / \texttt{wait} / \texttt{trigger}} —— \textbf{区域间通信}：轻量级的同步令牌，让并发区域能安全地通信。
\end{itemize}
